\scanpage{78}
4°) Cas $d'=$ ens des diagrammes finis, $d''=\varnothing$, pas de conditions d'exactitude.

Comme $d'$ est petit, les $R\in\mathrm{Ob}\,\Delta_0$ qui 2-représentent des th. algébriques sont les $R$ ess. petites (stables par $\varprojlim$ finies). \deleted{\unsure{Donnons-nous|Ramenons} \unsure{donc} $R$, lui \unsure{associons} $\unsure{\rho(R)}=\unsure{R^\circ}$ \unsure{Hom}\, $\unsure{(R,\mathrm{Ens})}$}

\[
\left\{\begin{aligned}
&T\quad\text{théorie algébrique pour }\Delta\\
&R\in\mathrm{Ob}\,\Delta\quad\text{qui représente }T\\
&S=T(\mathrm{Ens})=\underline{\mathrm{Hom}}_{\mathrm{lex}}(R,\mathrm{Ens}))\\
&\Sigma\subset S\quad\text{(structures représentables)}
\end{aligned}\right.
\]

Tout se détermine mutuellem\textsuperscript{t}

\[
\begin{aligned}
R:\qquad&T=T_R:C\longrightarrow\underline{\mathrm{Hom}}_{\mathrm{lex}}(R,C)\\
&S=\deleted{\unsure{T(\mathrm{E})}}\;\underline{\mathrm{Hom}}_{\mathrm{lex}}(R,\mathrm{Ens})\\
&\Sigma\cong R^\circ\subset S\quad\text{par foncteurs représ.}\\[2pt]
\Sigma:\qquad&\unsure{R|\ell}=\Sigma^\circ\\
&T:\unsure{R|C}\longmapsto\underline{\mathrm{Hom}}_{\mathrm{lex}}(\Sigma^\circ,C)\\
&S=\underline{\mathrm{Hom}}_{\deleted{\unsure{\mathrm{en}}}\mathrm{lex}}(S^\circ,C),\quad\text{contenant }\Sigma\text{ comme sous-catégorie pleine.}\\
&\deleted{\unsure{\text{ss-cat. repr.}}}
\end{aligned}
\]

\deleted{\unsure{Reste à voir une}} Autre façon de décrire $S$ de $R$ ou $\Sigma$

% slanted
\margin{cf SGA 4 I 8.7.5 où (i) $\unsure{\Sigma\in S|\Sigma_i\in S}$ où $\unsure{F|\Sigma}\in\underline{\mathrm{Hom}}(\unsure{\Sigma^\circ},\mathrm{Ens})$ (ii) $\Sigma$ strictement génératrice (iii) $\Sigma$ stable par petites $\varinjlim$, $\Sigma$ petite, le foncteur \unsure{:} ex. à dr.}

\uline{Lemme} \deleted{\unsure{On}} L'inclusion $\Sigma\longrightarrow S$ définit une \uline{équivalence}
\[
\mathrm{Ind}(\Sigma)\longrightarrow S
\]

\margin{SGA 4 I 8.9.5.}
[ En effet, \deleted{\unsure{est}} on sait que $\Sigma\longrightarrow S$ est universel pour des foncteurs $\Sigma\longrightarrow D$ ex. à dr, où $D$ a des $\varinjlim$ quelcon. Or montrons que $\Sigma\longrightarrow\mathrm{Ind}(\Sigma)$ est \deleted{\unsure{ident}} itou. \inserted{On sait que $\mathrm{Ind}(\Sigma)$ a des $\varinjlim$ finies (donc quelc.) car $\Sigma$ en a, et $\Sigma\longrightarrow\mathrm{Ind}(\Sigma)$ y comm.} \deleted{En effet} les foncteurs \uline{quelc.} $\Sigma\longrightarrow D$ \deleted{\unsure{sont les}}\inserted{correspondent}
\margin{SGA 4 I 8.7.3}
aux foncteurs $\mathrm{Ind}(\Sigma)\longrightarrow D$ qui commutent aux $\varinjlim$. Ceux qui commutent aux $\varinjlim$ finies (donc quelconques) sont ceux qui sur $\Sigma$ y commutent [nécessité triviale, pour la suffisance, on \deleted{note que \unsure{cela} revient aux \unsure{copr.} finies, initiaux, sommes, conoyaux,} \inserted{note que par I 8.9.8 $\mathrm{Ind}(\Sigma)\longrightarrow\mathrm{Ind}(D)$ \deleted{commute} est exact à droite,} or $\mathrm{Ind}\,D\longrightarrow D$ l'est aussi (c'est un adjoint à g, qui commute à \uline{ttes} limites inductives) ok.

\uline{Corollaire} \deleted{Pour tout} \inserted{Soit $X\in S$. Pour que} $X\in\Sigma$, \inserted{il f. et s. que} le foncteur $\mathrm{Hom}(X,-)$ sur $S$ commute aux $\varinjlim$ filtrantes.

% slanted
\margin{\uline{NB} Marche pour tte catégorie karoubienne ! cf aussi SGA 4 I 8.7.8}

C'est évidemment nécessaire, par construction des Hom dans $\mathrm{Ind}(S)$.

C'est suffisant, car \deleted{si $X\in S$,} On aura $X=\varinjlim_i X_i$ ($X_i\in\mathrm{Ob}\,\Sigma$), filtrante, donc le morphisme identique se factorise par un des $X_{\unsure{j|i}}$, donc on a un $X_i\longrightarrow X$ ayant une section, donc $X$ est une \deleted{\unsure{image}} image projecteur, or $\Sigma$ est \deleted{donc} stable \inserted{pour ça} (c'est un noyau) \deleted{et $\Sigma\longrightarrow S$, est exact à dr, donc} OK. $]$
% archive sheet number: 77
